\section{The Gaussian groups \(G_{29},G_{31}\): exact determinant obstruction}
\label{sec:g31-barrier}

The two natural rank-four Gaussian reflection representations do not
supply a dominant nonsplit Weil-sixfold family. This conclusion is based
on an exhaustive exact calculation from published matrices, followed by
D-DET-01; it is not a conclusion from a numerical period experiment.

\subsection*{D-G31-01. Published generators and exhaustive local determinants}

Put \(K=\mathbb Q(i)\), and use the standard Hermitian form on \(K^4\).
Davis \cite[Theorem 4.1, proof, p.~307]{G34:Davis} gives the four
generators of \(G_{29}\):
\[
s_1=\frac12\begin{pmatrix}
1&-1&-1&-1\\-1&1&-1&-1\\-1&-1&1&-1\\-1&-1&-1&1
\end{pmatrix},\quad
s_2=\begin{pmatrix}0&i&0&0\\-i&0&0&0\\0&0&1&0\\0&0&0&1\end{pmatrix},
\]
\[
s_3=\begin{pmatrix}0&1&0&0\\1&0&0&0\\0&0&1&0\\0&0&0&1\end{pmatrix},\qquad
s_4=\begin{pmatrix}1&0&0&0\\0&0&1&0\\0&1&0&0\\0&0&0&1\end{pmatrix}.
\]
His Theorem 4.16, p.~311, adds
\(s_5=\operatorname{diag}(1,1,-1,1)\) to generate \(G_{31}\).
All five are self-adjoint involutive reflections.

For an element \(g\), let \(U_g=\operatorname{im}(1-g)\),
\(k(g)=\dim_KU_g\), and choose a \(K\)-basis of \(U_g\). Define
\[
b(g)=\det(1-g|_{U_g}),\qquad
q_+(g)=\frac{\det(h|_{U_g})}{|b(g)|}.
\]
Since all group determinants are \(\pm1\),
\(\overline{b(g)}=(-1)^{k(g)}\det(g)^{-1}b(g)\), so \(b(g)\)
is rational or purely imaginary rational. Hence \(q_+(g)\) is a
positive rational number. Its class modulo Gaussian norms is independent
of the chosen moving-space basis and of conjugation in the group.

The complete exact enumeration gives:
\[
\begin{array}{c|r|r|r|r}
&|G|&\text{conjugacy classes}&\text{reflections}&
 \text{nonidentity involutions}\\\hline
G_{29}&7680&37&40&231\\
G_{31}&46080&59&60&511
\end{array}
\]
For \(G_{29}\), every local \(q_+\) is a Gaussian norm. For
\(G_{31}\), every class of moving rank at most three, and every
involution class, has norm \(q_+\). The three exceptions are:
\[
\begin{array}{c|r|r|c|c|c|c}
\text{certificate index}&\text{class size}&\text{order}&k(g)&b(g)&q_+(g)
 &\operatorname{age}(g)-k(g)/2\\\hline
41&960&12&4&3i&3&1/2\\
43&960&12&4&-3i&3&-1/2\\
46&960&6&4&3&3&0
\end{array}
\]
The indices are those of our enumeration, not imported character-table
labels. In the selected bases all three moving Gram determinants equal
nine. Their nonzero eigenvalue exponents are respectively
\((4,7,8,11)/12\), \((1,4,5,8)/12\), and \((1,2,4,5)/6\).

Here is a transparent specification of the exhaustive certificate.
\texttt{enumerate\_gaussian\_groups.cpp} stores twice each matrix as
32 signed integers, the real and imaginary entries in row order. Every
product divides its integer numerator by two, with exact divisibility
checked. Breadth-first multiplication by the four or five displayed
generators enumerates the entire generated group. All matrices satisfy
\(g^*g=I\) exactly. Breadth-first conjugation by the same generators
partitions it into complete conjugacy classes; the class sizes sum to
the full enumerated order. Thus the class list is exhaustive.

\texttt{check\_local\_discriminants.py} independently performs rational
Gaussian elimination on each class representative. It takes independent
columns of \(1-g\), computes the induced moving-space matrix, its Gram
determinant, and \(b(g)\). It tests the Gaussian norm criterion by
factoring the positive rational number \(q_+\): each prime congruent
to three modulo four must have even valuation. Eigenvalue multiplicities
are computed by the exact finite Fourier formula in
\(\mathbb Q[\zeta_{\operatorname{lcm}(4,\operatorname{ord}g)}]\), reducing
polynomials by the cyclotomic polynomial, not by floating-point roots.
Every multiplicity is checked to be a nonnegative integer.

The files \texttt{G29\_classes.json}, \texttt{G31\_classes.json}, and
\texttt{G29\_local\_discriminants.json},
\texttt{G31\_local\_discriminants.json} retain every representative,
class size, order, local matrix determinant, Gram determinant, spectrum
and norm test. Together with the short source code, these constitute
the full finite computation certificate used here.

\subsection*{D-G31-02. The universal nonsplit-source obstruction}

Consider a genus-zero product-one generating tuple with balanced
parabolic multiplicity space. D-DET-01 gives its signed determinant
class as
\[
                  \left[\prod_j\frac{\det(h|_{U_{g_j}})}{b(g_j)}\right].
\]
The displayed representative is positive: it differs from the determinant
of the anti-Hermitian parabolic cup matrix by a positive norm, and for
a balanced Hermitian space that anti-Hermitian determinant is positive.
Since each factor is real or purely imaginary, the positive representative
equals the product of its absolute values. Therefore its norm class is
\[
                              \left[\prod_jq_+(g_j)\right].
\]
For \(G_{29}\) this is always a norm. Thus a balanced genus-zero
\(G_{29}\) source is necessarily split.

For \(G_{31}\), a nonsplit balanced source must contain an odd number
of the three exceptional classes, in particular at least one inertia
element of moving rank four. A sixfold has total moving rank
\(D=2\cdot4+6=14\). If its branch count is \(r\), that rank-four
entry and the other nonidentity entries give \(D\ge4+(r-1)\), hence
\[
                              r\le11.
\]
Its Hurwitz dimension and period-image dimension are consequently at most
eight. This is a necessary-dimension obstruction, not a claim that the
upper bound is attained by every such tuple.

Allowing higher-genus bases does not restore dominance. For a base of
genus \(\gamma\), the same sixfold condition is
\(D=14-8\gamma\). At genus one there are at most six branch points
and at most six parameters; at larger genus the dimension equation is
impossible. Therefore neither natural representation gives a dominant
nonsplit sixfold curve source, for arbitrary initial inertia orders,
Nielsen class, ordering, or base genus.

\subsection*{D-G31-03. An explicit nonsplit tuple and its rank-two boundary}

The determinant obstruction is visible in a concrete connected family.
Besides the five published generators define
\[
u=\frac12\begin{pmatrix}
1&-i&-1&i\\i&1&i&1\\-1&-i&1&i\\-i&1&-i&1
\end{pmatrix},\qquad
c=\frac12\begin{pmatrix}
0&1-i&-1-i&0\\-1+i&0&0&1+i\\
1-i&0&0&1+i\\0&-1+i&-1-i&0
\end{pmatrix}.
\]
Exact multiplication gives
\[
          c=(s_1s_4s_5s_3s_2u)^{-1},\qquad
          c^6=1,\quad \det(1-c)=3.
\]
The matrix \(u\) is a reflection in the conjugacy orbit of \(s_5\).
The eigenvalues of \(c\) have exponents \((1,2,4,5)/6\), so its
age is two and its moving rank is four. The eleven-entry tuple
\[
       (s_1,s_1,s_2,s_2,s_1,s_4,s_5,s_3,s_2,u,c)
\]
has product one and generates \(G_{31}\), since its tail includes
all five published generators. Riemann existence gives connected covers;
their branch orders are \((2^{10},6)\) and their genus is
\[
           1+\frac{46080}{2}\left(-2+10\cdot\frac12+\frac56\right)
                       =88321.
\]
Chevalley--Weil gives signature \((3,3)\): each reflection contributes
age \(1/2\), and \(c\) and its inverse each contribute age two.
The exact parabolic quotient has relation and coboundary ranks four,
dimension six, and anti-Hermitian cup determinant \(3/8\). Hence its
signed discriminant has Gaussian norm class three, which is nonsplit.
The source has dimension eight; its exact period rank has not been
computed and is not asserted to be eight.

There is a useful degeneration of precisely the form sought for a
nonsplit sixfold. Bubble off the first two equal pairs separately.
The two new base nodes have trivial monodromy. The central normalization
has branch tuple
\((s_1,s_4,s_5,s_3,s_2,u,c)\), still generating the full group;
its natural multiplicity space has dimension two, signature \((1,1)\),
and exact cup determinant \(3/2\). It is an anisotropic Hermitian plane,
since three is not a Gaussian norm. The bubble components are rational
cyclic double covers.

The dual graph has one central vertex, bubble vertices
\(G/\langle s_1\rangle\) and \(G/\langle s_2\rangle\), and two
copies of \(G\) as edge sets. Its natural-representation graph
multiplicity is
\[
             2\dim V-\dim V^{s_1}-\dim V^{s_2}=8-3-3=2.
\]
The vanishing circles satisfy exactly the vertex-incidence relations;
this follows by Mayer--Vietoris on the punctured normalization
components. A positive smoothing applies a positive Dehn twist at
each node. The sum of their rank-one transvection forms has image the
full span of those circles, since its positive quadratic form is
nondegenerate on that span. Taking the exact natural multiplicity
functor gives a logarithm of rank two. The weight-one part of the
normalization is the anisotropic plane just computed. Thus the
degeneration has graph rank two and an anisotropic residual Weil surface,
with no fictitious maximal isotropic rank-three limit.

\texttt{find\_nonsplit\_tuple.py} reconstructs the reflection orbit,
verifies the product, generation witness, exact spectrum and both
parabolic Gram matrices, and writes \texttt{nonsplit\_tuple.json}.
This establishes the source and its determinant, not algebraicity of
its sixfold Weil classes: the higher-order inertia requires a new
supplier/character-coverage argument, and the source is too small to
dominate the nine-dimensional nonsplit moduli space.
